% !PRINTED 636
\weberpage{636}
%<<CONT
and if one multiplies by $b_0^{e-1}$ and sets $b_0\eta = \omega$, one obtains an equation
of the form (1).

5.\ \gesperrt{A function $\omega$ of the field $\Omega$ is an integral function if it
satisfies any equation}

\weq{5}{\psi(\omega) = \omega^m + c_1\omega^{m-1} + \cdots + c_{m-1}\omega + c_m = 0}

\gesperrt{in which $c_1, c_2, \ldots, c_m$ are integral rational functions, even if this is
not the equation of lowest degree.}

For if $\varphi(\omega) = 0$ is the equation of lowest degree for $\omega$, then $\psi(t)$
must, for a variable $t$, be divisible by $\varphi(t)$. That is simply the content of
equation (5). We have therefore

\[
\psi(t) = \chi(t)\varphi(t),
\]

in which $\psi(t)$ is likewise an integral function of $t$.
% [sic] printed psi(t), error for chi(t): the clause describes the factor chi(t) just introduced in psi(t) = chi(t) phi(t), and psi(t) is already integral by hypothesis
If now, as was assumed,
$\psi(t)$ has integral coefficients, then by a theorem of \gesperrt{Gauss} the
coefficients of $\varphi(t)$ and of $\chi(t)$ must also be integral functions. Cf.\ Bd.\ I,
\S\,2, \S\,20, \S\,148\,$^1$).

From this there follows the principal theorem below:

6.\ \gesperrt{Sum, difference, product of two integral functions in $\Omega$ are again
integral functions.}

For if $\omega'$, $\omega''$ are two integral functions in $\Omega$ which satisfy the
equations:

\weq{6}{
\begin{aligned}
\omega'^{\,n'} + b_1'\omega'^{\,n'-1} + \cdots + b_{n'-1}'\omega' + b_{n'}' &= 0, \\
\omega''^{\,n''} + b_1''\omega''^{\,n''-1} + \cdots + b_{n''-1}''\omega'' + b_{n''}'' &= 0,
\end{aligned}
}

then let

\[
\omega_1, \omega_2, \ldots, \omega_m
\]

denote the products

\[
\omega'^{\,\nu'}\omega''^{\,\nu''}; \quad \nu' = 0, 1, \ldots, n'-1, \quad
\nu'' = 0, 1, \ldots, n''-1,
\]

and let $\omega$ denote one of the three functions $\omega' \pm \omega''$,
$\omega'\omega''$; then with the help of equations (6) one may set:

\weq{7}{
\begin{aligned}
\omega\omega_1 &= x_{1,1}\omega_1 + \cdots + x_{1,m}\omega_m, \\
&\cdots\cdots\cdots\cdots\cdots\cdots\cdots \\
\omega\omega_m &= x_{m,1}\omega_1 + \cdots + x_{m,m}\omega_m,
\end{aligned}
}

%CONT>>
\begin{quote}
$^1$) In this proof the decomposability of integral rational functions into linear factors
need not be assumed, but only the decomposability into irreducible factors. One may
therefore still stipulate an arbitrary domain of rationality for the numerical
coefficients. In what follows, however, we shall nevertheless have to assume the
decomposability of an integral rational function of $z$ into linear factors, i.\ e.\ the
fundamental theorem of algebra.
\end{quote}
